\documentclass[12pt]{article}
\usepackage[T1]{fontenc}
\usepackage[utf8]{inputenc}
\usepackage{lmodern}
\usepackage{textcomp}
\usepackage{enumitem}
\usepackage{geometry}
\usepackage{graphicx}
\usepackage{amsmath, amssymb}
\geometry{margin=1in}
\begin{document}
\begin{center}
Chemistry - Undergraduate Level Assessment 20 \
Total Marks: 40 \
Answer all questions.
\end{center}

\section*{Multiple Choice (10 marks)}
\begin{enumerate}[label=\arabic*.]
\item Which of the following compounds has a 1 H resonance approximately 1.55 kHz away from TMS on a spectrometer with a 12.0 T magnet?
\begin{enumerate}[label=(\alph*)]
    \item CH 3 F
    \item CH 3 Cl
    \item CH 3 Br
    \item CH 3 I
\end{enumerate}
\item Calculate the Q-factor for an X-band EPR cavity with a resonator bandwidth of 1.58 MHz.
\begin{enumerate}[label=(\alph*)]
    \item Q = 1012
    \item Q = 2012
    \item Q = 3012
    \item Q = 6012
\end{enumerate}
\item Which of the following is required for both paramagnetism and ferromagnetism?
\begin{enumerate}[label=(\alph*)]
    \item Strong oxidizing conditions
    \item Low-spin electron configuration
    \item Metallic physical properties
    \item Unpaired electrons
\end{enumerate}
\item Which nuclide has an NMR frequency of 115.5 MHz in a 20.0 T magnetic field?
\begin{enumerate}[label=(\alph*)]
    \item 17 O
    \item 19 F
    \item 29 Si
    \item 31 P
\end{enumerate}
\item Cobalt-60 is used in the radiation therapy of cancer and can be produced by bombardment of cobalt-59 with which of the following?
\begin{enumerate}[label=(\alph*)]
    \item Neutrons
    \item Alpha particles
    \item Beta particles
    \item X-rays
\end{enumerate}
\end{enumerate}

\section*{True / False (10 marks)}
\textit{Answer True or False}
\begin{enumerate}[label=\arabic*.]
\setcounter{enumi}{5}
\item True or False: T$_{1}$, unlike T$_{2}$, is sensitive to molecular motions at the Larmor frequency.
\item True or False: The total entropy of the system plus surroundings increases during a spontaneous process.
\item True or False: Planck's quantum theory explains the emission spectra of diatomic molecules in terms of continuous energy distribution.
\item True or False: The hydrogen atom in CHCl 3 does not undergo rapid intermolecular exchange, resulting in multiple peaks in its 1 H NMR spectrum.
\item True or False: The limiting high-temperature molar heat capacity at constant volume (C\_V) of a gas-phase diatomic molecule is 2.5 R.
\end{enumerate}

\section*{Long Form Response (20 marks)}
\textit{Answer each question with a clear and complete explanation.}
\begin{enumerate}[label=\arabic*.]
\setcounter{enumi}{10}
\item Discuss the factors influencing the thermal stability of hydrides in group-14 elements and analyze their stability order, providing a detailed explanation of each compound's characteristics.
\item Discuss the significance of electromagnetic radiation absorption by proteins at a wavelength of 190 nm, particularly focusing on its implications in the ultraviolet region of the spectrum.
\end{enumerate}

\vfill
\noindent\textit{End of Paper}
\end{document}